% Book 3, Year 1, chapter 6 -- strings, structures and files.
% Provenance (never printed): sources/book3/BRIEFS.md, chapter 6.
% ledger: b1:ascii-nul, b1:snprintf, b1:strtol, b1:psabi-struct-layout, b1:c-fopen-modes, b1:x86-little-endian
\chapter{Strings, Structures and Files}\label{ch:b1:c-strings-structs-files}

A program keeps a person's name in a ten-byte field and their age just after it. Someone types a twelve-letter name;
the program copies it; the name prints correctly, and the age has become 0. Nothing crashed, no warning appeared, and
the next program to read the age will believe it. This chapter explains the three facts behind that story: how C
stores text, how it lays out a group of \omterm{def:b1:c-strings-structs-files:struct}{fields} in memory, and how those bytes reach a file. With them, a C program
can read the town's weather file robustly and write it back in a form that a program can reach record by record.

\begin{recall}
A \omterm{def:b1:c-strings-structs-files:text-binary}{text file} is a sequence of characters cut into lines, and a file of records holds one record per line, its \omterm{def:b1:c-strings-structs-files:struct}{fields}
separated by a character such as a comma. Each character is stored as a number, its code: ASCII for the letters of
English, UTF-8 for every script, ASCII characters keeping their one-byte codes. Programs read standard input and write
standard output.
\end{recall}

\section{Null-terminated strings}\label{sec:b1:c-strings-structs-files:strings}

C has no string type. It has \omterm{def:b1:pointers-and-arrays:array}{arrays} of \code{char}, and a convention.

\begin{definition}[Null-terminated string]\label{def:b1:c-strings-structs-files:string}
In C, a \emph{null-terminated string}\index{null-terminated string}\index{string (C)} is a sequence of characters
stored in consecutive \code{char} and followed by a byte of value 0, the null character \texttt{'\textbackslash{}0'}
(code 0, NUL, in ASCII). Its length is the number of characters before that byte. A string literal \code{"rain"} is an
\omterm{def:b1:pointers-and-arrays:array}{array} of 5 \code{char}, the last one zero, that the program must not modify.
\end{definition}

\omcode{code/bachelor-1/06-c-strings-structs-files/listings/strlen_demo.c}{2}{16}{The characters of a string, and
its length}\label{lst:b1:c-strings-structs-files:strlen}

\omoutput{code/bachelor-1/06-c-strings-structs-files/outputs/strlen_demo.txt}{Eight codes, then the zero byte; cutting
the string takes one assignment.}

\code{strlen} counts the characters until it meets the zero byte: it reads the whole string, so its cost is linear in
the length, and calling it in the condition of a loop over the string, as \cref{lst:b1:c-strings-structs-files:strlen}
does for brevity, makes the loop quadratic. \code{sizeof} gives the size of the \omterm{def:b1:pointers-and-arrays:array}{array}, fixed at compile time and
unrelated to the length of the string it currently holds. Every function of \code{<string.h>} relies on the zero byte:
\code{strcpy(dst, src)} copies the characters of \code{src} and its zero into \code{dst}, \code{strcat(dst, src)}
copies them after the end of \code{dst}, \code{strcmp} compares two strings character by character. None of them knows
the size of \code{dst}.

\begin{definition}[Buffer overflow]\label{def:b1:c-strings-structs-files:buffer-overflow}
A \emph{buffer overflow}\index{buffer overflow} is a write past the end (or before the start) of an \omterm{def:b1:pointers-and-arrays:array}{array}, into the
memory that follows it. It is \omterm{def:b1:c-types-and-control:ub}{undefined behaviour}; in practice, it changes whatever the program keeps there.
\end{definition}

Here is the story of the introduction, built three ways: plainly, with the address sanitizer, and with the
optimisation that this machine's \omterm{def:b1:compilation-pipeline:compiler}{compiler} accompanies with extra checks.

\omcode{code/bachelor-1/06-c-strings-structs-files/listings/overflow.c}{2}{17}{Twelve letters for a ten-byte
field}\label{lst:b1:c-strings-structs-files:overflow}

\omoutput{code/bachelor-1/06-c-strings-structs-files/outputs/overflow.txt}{Silent corruption, a sanitizer that sees
nothing, the same copy into a plain \omterm{def:b1:pointers-and-arrays:array}{array} caught, and the optimised build stopped.}\label{out:b1:c-strings-structs-files:overflow}

\begin{omfigure}
\begin{tikzpicture}[b/.style={draw, minimum width=7mm, minimum height=7mm, font=\small\ttfamily, inner sep=0pt}]
\foreach \c [count=\i from 0] in {C,h,r,i,s,t,a,b,e,l} {\node[b, fill=omCode!8] at (\i*0.7,0) {\c};}
\node[b, fill=omThm!20] at (10*0.7,0) {l};
\node[b, fill=omThm!20] at (11*0.7,0) {e};
\node[b, fill=omThm!20] at (12*0.7,0) {0};
\node[b, fill=omThm!20] at (13*0.7,0) {0};
\node[b, fill=omThm!20] at (14*0.7,0) {0};
\node[b, fill=omThm!20] at (15*0.7,0) {0};
\draw[decorate, decoration={brace, amplitude=4pt}] (-0.35,0.45) -- (6.65,0.45) node[midway, above=4pt, font=\small]
  {\code{name}, 10 bytes};
\draw[decorate, decoration={brace, amplitude=4pt}] (6.65,0.45) -- (8.05,0.45) node[midway, above=4pt, font=\small]
  {padding};
\draw[decorate, decoration={brace, amplitude=4pt}] (8.05,0.45) -- (10.85,0.45) node[midway, above=4pt, font=\small]
  {\code{age}};
\node[font=\scriptsize, align=center] at (9.45,-0.85) {was \code{13 00 00 00} (19),\\now \code{00 00 00 00}};
\node[font=\scriptsize, align=left, anchor=west] at (-0.35,-0.85) {offsets 0 to 15 of \code{struct person}};
\end{tikzpicture}
\par\smallskip{\small The thirteen bytes written by \code{strcpy}. The last two letters fall into the two unused bytes
that keep \code{age} at an offset multiple of 4; the zero byte that ends the string lands on the low byte of
\code{age}, stored first, and 19 becomes 0.}
\end{omfigure}

The address sanitizer protects the space around each variable, not the inside of a \omterm{def:b1:c-strings-structs-files:struct}{structure}: the write stays within
the 16 bytes of \code{p}, so it is not reported. Into a plain \omterm{def:b1:pointers-and-arrays:array}{array} of 10 bytes the same copy is caught at once. The
third build shows a defence many systems add: when optimising, this \omterm{def:b1:compilation-pipeline:compiler}{compiler} defines \code{_FORTIFY_SOURCE}, which
replaces \code{strcpy} by a version that knows the size of \code{p.name} when the \omterm{def:b1:compilation-pipeline:compiler}{compiler} does, and stops the program
on an overflow. None of the three should be relied on. The size must travel with the buffer, and every copy must
respect it.

\omcode{code/bachelor-1/06-c-strings-structs-files/listings/safe_copy.c}{2}{16}{A bounded copy that reports
truncation}\label{lst:b1:c-strings-structs-files:safe}

\omoutput{code/bachelor-1/06-c-strings-structs-files/outputs/safe_copy.txt}{\code{snprintf} keeps 9 characters and the
zero byte, and returns the length it would have needed.}

\code{snprintf(buf, size, format, ...)} writes at most \code{size} bytes, the last always zero, and returns the length
the full result would have had: a result greater than or equal to \code{size} means the text was cut. The program then
decides: refuse, warn, or allocate a bigger buffer.

A program receives its command-line arguments as strings: \code{main(int argc, char *argv[])} gets their number and an
\omterm{def:b1:pointers-and-arrays:array}{array} of \omterm{def:b1:pointers-and-arrays:pointer}{pointers} to them, \code{argv[0]} being the program's name. Turning them into numbers is a small lesson in
robustness.

\omcode{code/bachelor-1/06-c-strings-structs-files/listings/args.c}{3}{21}{Two conversions of the same
arguments}\label{lst:b1:c-strings-structs-files:args}

\omoutput{code/bachelor-1/06-c-strings-structs-files/outputs/args.txt}{\code{atoi} cannot say that something went
wrong.}\label{out:b1:c-strings-structs-files:args}

\code{atoi} returns 0 for \code{hello}, as for the string \code{0}, and 7 for \code{7x}; on a number too large for an
\code{int} its behaviour is undefined, here a value cut to 32 bits. \code{strtol} reports, through its second argument,
where the digits stopped, and through \code{errno}, a variable of the C library that it sets to \code{ERANGE}, an
overflow of \code{long}. The last argument fits in a \code{long} but not in an \code{int}: a caller wanting an
\code{int} still has to check the range (\cref{exo:b1:c-strings-structs-files:12}).

\section{Structures, unions and enumerations}\label{sec:b1:c-strings-structs-files:structs}

\begin{definition}[Structure, field]\label{def:b1:c-strings-structs-files:struct}
A \emph{structure}\index{structure (C struct)} in C is a type made of several named members, the \emph{fields}\index{field
(of a structure)}, of possibly different types, stored in one block of memory in the order of their declaration. A
field is reached with \code{.} from a structure and with \code{->} from a \omterm{def:b1:pointers-and-arrays:pointer}{pointer} to one: \code{p->x} means
\code{(*p).x}.
\end{definition}

\omcode{code/bachelor-1/06-c-strings-structs-files/listings/point.c}{2}{32}{Initialising, reaching and copying
\omterm{def:b1:c-strings-structs-files:struct}{structures}}\label{lst:b1:c-strings-structs-files:point}

\omoutput{code/bachelor-1/06-c-strings-structs-files/outputs/point.txt}{A \omterm{def:b1:c-strings-structs-files:struct}{structure} is copied as a whole, \omterm{def:b1:c-functions-and-stack:by-value}{by value}, like
an \code{int}.}

\code{typedef} gives a type a second name, so that \code{point} can be written instead of a full \omterm{def:b1:c-strings-structs-files:struct}{structure} type.
\emph{Designated initialisers}, \code{.name = ...}, set \omterm{def:b1:c-strings-structs-files:struct}{fields} by name in any order and set every other field to zero.
A \omterm{def:b1:c-strings-structs-files:struct}{structure} passed to a function is copied, like any argument (\cref{ch:b1:c-functions-and-stack}); assignment copies
every field, including \omterm{def:b1:pointers-and-arrays:array}{arrays} inside it, which an \omterm{def:b1:pointers-and-arrays:array}{array} alone cannot do. To let a function change a \omterm{def:b1:c-strings-structs-files:struct}{structure}, or to
avoid copying a large one, pass its address.

\begin{definition}[Enumeration]\label{def:b1:c-strings-structs-files:enum}
An \emph{enumeration}\index{enumeration (C)} declares named integer constants: \code{enum kind { CIRCLE, RECTANGLE };}
makes \code{CIRCLE} 0 and \code{RECTANGLE} 1, and \code{enum kind} a type whose values are meant to be those names.
\end{definition}

\begin{definition}[Union, tagged union]\label{def:b1:c-strings-structs-files:union}
A \emph{union}\index{union (C)} is a type whose members all start at the same address and share the same storage: its
size is that of its largest member, and only the member last written holds a meaningful value. A \emph{tagged
union}\index{tagged union} is a \omterm{def:b1:c-strings-structs-files:struct}{structure} holding a union and a field, the tag, that records which member is in use.
\end{definition}

\omcode{code/bachelor-1/06-c-strings-structs-files/listings/shapes.c}{2}{38}{A shape is a circle or a
rectangle}\label{lst:b1:c-strings-structs-files:shapes}

\omoutput{code/bachelor-1/06-c-strings-structs-files/outputs/shapes.txt}{The \omterm{def:b1:c-strings-structs-files:union}{union} takes the size of its larger
member; a lying tag gives a wrong answer and no error.}

\begin{omfigure}
\begin{tikzpicture}[b/.style={draw, minimum height=7mm, font=\small, inner sep=1pt, anchor=west}]
\foreach \y/\lab in {0/circle, -1.2/rectangle} {
  \node[b, minimum width=12mm, fill=omCode!8] at (0,\y) {\code{kind}};
  \node[b, minimum width=12mm, fill=black!8] at (1.2,\y) {pad};
  \node[font=\small, anchor=east] at (-0.2,\y) {\lab};
}
\node[b, minimum width=24mm, fill=omDef!12] at (2.4,0) {\code{r}};
\node[b, minimum width=24mm, fill=black!4] at (4.8,0) {unused};
\node[b, minimum width=24mm, fill=omDef!12] at (2.4,-1.2) {\code{w}};
\node[b, minimum width=24mm, fill=omDef!12] at (4.8,-1.2) {\code{h}};
\foreach \x/\o in {0/0, 1.2/4, 2.4/8, 4.8/16, 7.2/24} {\node[font=\scriptsize] at (\x,0.6) {\o};}
\draw[decorate, decoration={brace, amplitude=4pt, mirror}] (2.4,-1.7) -- (7.2,-1.7) node[midway, below=4pt,
  font=\small] {the union \code{as}: 16 bytes, read through the member the tag names};
\end{tikzpicture}
\par\smallskip{\small A \omterm{def:b1:c-strings-structs-files:union}{tagged union} of 24 bytes. The same 16 bytes hold a radius or a width and a height; the tag
(4 bytes, then 4 bytes of \omterm{def:b1:c-strings-structs-files:alignment}{padding} so the \code{double} starts at offset 8) says which reading is meaningful.}
\end{omfigure}

The function \code{area} must check the tag before reading the \omterm{def:b1:c-strings-structs-files:union}{union}, and the \code{switch} without a
\code{default} lets \code{-Wall} warn when a new kind is added to the \omterm{def:b1:c-strings-structs-files:enum}{enumeration} but not to the \code{switch}. Nothing
in C prevents the tag and the contents from disagreeing; the Year~2 volume meets languages whose sum types make that
disagreement impossible.

\section{Layout and alignment}\label{sec:b1:c-strings-structs-files:layout}

The \omterm{def:b1:c-strings-structs-files:struct}{structure} of the opening story occupies 16 bytes for 14 bytes of \omterm{def:b1:c-strings-structs-files:struct}{fields}. The two extra bytes come from a
requirement of the processor and its calling conventions.

\begin{definition}[Alignment, padding]\label{def:b1:c-strings-structs-files:alignment}
The \emph{alignment}\index{alignment} of a type is a power of two by which the address of every object of that type
must be divisible; on this machine it equals the size for \code{char}, \code{short}, \code{int}, \code{long},
\code{double} and \omterm{def:b1:pointers-and-arrays:pointer}{pointers}. \emph{Padding}\index{padding} is the unused bytes that the \omterm{def:b1:compilation-pipeline:compiler}{compiler} inserts in a \omterm{def:b1:c-strings-structs-files:struct}{structure}
so that each field lies at an offset multiple of its alignment, and at the end so that the size is a multiple of the
largest alignment of a field, which keeps every element of an \omterm{def:b1:pointers-and-arrays:array}{array} of \omterm{def:b1:c-strings-structs-files:struct}{structures} aligned.
\end{definition}

\omcode{code/bachelor-1/06-c-strings-structs-files/listings/layout.c}{2}{33}{The same four \omterm{def:b1:c-strings-structs-files:struct}{fields} in two
orders}\label{lst:b1:c-strings-structs-files:layout}

\omoutput{code/bachelor-1/06-c-strings-structs-files/outputs/layout.txt}{24 bytes, then 16, for the same 14 bytes of
\omterm{def:b1:c-strings-structs-files:struct}{fields}.}\label{out:b1:c-strings-structs-files:layout}

\begin{omfigure}
\begin{tikzpicture}[x=5.2mm, y=1mm]
\newcommand{\fld}[5]{\draw[fill=#4] (#1,#2) rectangle (#1+#3,#2+6); \node[font=\scriptsize] at (#1+#3/2,#2+3) {#5};}
\foreach \o in {0,4,...,24} {\node[font=\scriptsize] at (\o,8.5) {\o};}
\node[font=\small, anchor=east] at (-0.4,3) {\code{loose}};
\fld{0}{0}{1}{omCode!8}{s} \fld{1}{0}{7}{black!12}{padding} \fld{8}{0}{8}{omDef!12}{\code{temp}}
\fld{16}{0}{1}{omCode!8}{f} \fld{17}{0}{3}{black!12}{pad} \fld{20}{0}{4}{omDef!12}{\code{date}}
\node[font=\small, anchor=east] at (-0.4,-9) {\code{tight}};
\fld{0}{-12}{8}{omDef!12}{\code{temp}} \fld{8}{-12}{4}{omDef!12}{\code{date}} \fld{12}{-12}{1}{omCode!8}{s}
\fld{13}{-12}{1}{omCode!8}{f} \fld{14}{-12}{2}{black!12}{pad}
\foreach \k in {0,...,23} {\draw[black!30] (\k,0) -- (\k,1);}
\end{tikzpicture}
\par\smallskip{\small The two layouts of \cref{lst:b1:c-strings-structs-files:layout}, with offsets in bytes (s:
\code{station}, f: \code{flag}). In \code{loose}, the \code{double} forces seven bytes of \omterm{def:b1:c-strings-structs-files:alignment}{padding} after the first
\code{char}; sorted by decreasing \omterm{def:b1:c-strings-structs-files:alignment}{alignment}, the \omterm{def:b1:c-strings-structs-files:struct}{fields} leave only two bytes at the end.}
\end{omfigure}

\begin{method}[Ordering the fields of a structure]\label{met:b1:c-strings-structs-files:order}
Declare the \omterm{def:b1:c-strings-structs-files:struct}{fields} by decreasing \omterm{def:b1:c-strings-structs-files:alignment}{alignment}: \code{double}, \omterm{def:b1:pointers-and-arrays:pointer}{pointers} and \code{long} first, then \code{int} and
\code{float}, then \code{short}, then \code{char} and \omterm{def:b1:pointers-and-arrays:array}{arrays} of \code{char}. No \omterm{def:b1:c-strings-structs-files:alignment}{padding} is then needed between \omterm{def:b1:c-strings-structs-files:struct}{fields},
and the size is the sum of the \omterm{def:b1:c-strings-structs-files:struct}{fields} rounded up to the largest \omterm{def:b1:c-strings-structs-files:alignment}{alignment}, the smallest size possible
(\cref{exo:b1:c-strings-structs-files:9}). Use \code{offsetof} and \code{sizeof} to check, never a count by hand.
\end{method}

The C standard leaves layout to the implementation; on this machine, the rules are those of the calling conventions
of \cref{ch:b1:c-functions-and-stack}, which every \omterm{def:b1:compilation-pipeline:compiler}{compiler} for it follows, so that a \omterm{def:b1:c-strings-structs-files:struct}{structure} built by one
library is read correctly by another.

\section{Text files and binary files}\label{sec:b1:c-strings-structs-files:files}

\begin{definition}[Stream, standard streams]\label{def:b1:c-strings-structs-files:stream}
A \emph{stream}\index{stream} is the C library's access to a file or a device: a \code{FILE *} returned by
\code{fopen}, through which the program reads or writes bytes in order, the library collecting them in a buffer and
passing them to the system in large blocks. Three \emph{standard streams}\index{standard streams} are open when
\code{main} starts: \code{stdin} (standard input), \code{stdout} (standard output) and \code{stderr} (standard error,
for messages, not fully buffered).
\end{definition}

\code{fopen(path, mode)} opens a \omterm{def:b1:c-strings-structs-files:stream}{stream}, or returns \code{NULL} and sets \code{errno}; \code{perror} prints the
reason. The mode \code{"r"} reads; \code{"w"} creates the file or empties it; \code{"a"} writes at the end; a
\code{"+"} allows both reading and writing; a \code{"b"} asks for binary mode, which changes nothing on this system but
matters on systems that translate line ends. \code{fclose} writes what is left in the buffer and can fail, so its
result is checked when writing.

\begin{definition}[Text file, binary file]\label{def:b1:c-strings-structs-files:text-binary}
A \emph{text file}\index{text file} holds characters only, usually in lines, so that a person can read it and any
program can parse it; a number in it is a sequence of digit characters. A \emph{binary file}\index{binary file} holds
bytes in a format defined by a program, such as the bytes of values copied from memory; it is smaller and faster to
read, but meaningless without its format.
\end{definition}

\omcode{code/bachelor-1/06-c-strings-structs-files/listings/write_text.c}{6}{21}{Three numbers as
text}\label{lst:b1:c-strings-structs-files:wtext}

\omcode{code/bachelor-1/06-c-strings-structs-files/listings/write_binary.c}{6}{21}{The same numbers as
bytes}\label{lst:b1:c-strings-structs-files:wbin}

\omoutput{code/bachelor-1/06-c-strings-structs-files/outputs/bytes.txt}{Both files, their sizes, and their bytes in
hexadecimal (\code{od -t x1}), then the binary one read back as 4-byte integers.}\label{out:b1:c-strings-structs-files:bytes}

\begin{omfigure}
\begin{tikzpicture}[b/.style={draw, minimum width=6.5mm, minimum height=6.5mm, font=\scriptsize\ttfamily,
                    inner sep=0pt}]
\node[font=\small, anchor=east] at (-0.5,0) {text, 17 bytes};
\foreach \c/\h [count=\i from 0] in {7/37,{\textbackslash n}/0a,1/31,0/30,0/30,0/30,{\textbackslash n}/0a,3/33,0/30,
  5/35,4/34,1/31,9/39,8/38,9/39,6/36,{\textbackslash n}/0a}
  {\node[b, fill=omCode!8] at (\i*0.65,0) {\h}; \node[font=\scriptsize] at (\i*0.65,0.55) {\c};}
\node[font=\small, anchor=east] at (-0.5,-1.4) {binary, 12 bytes};
\foreach \h [count=\i from 0] in {07,00,00,00,e8,03,00,00,78,56,34,12}
  {\node[b, fill=omDef!12] at (\i*0.65,-1.4) {\h};}
\foreach \a/\b/\v in {0/3/7, 4/7/1000, 8/11/305419896}
  {\draw[decorate, decoration={brace, amplitude=3pt, mirror}] (\a*0.65-0.3,-1.8) -- (\b*0.65+0.3,-1.8)
     node[midway, below=3pt, font=\scriptsize] {\v};}
\end{tikzpicture}
\par\smallskip{\small The same three numbers. In the \omterm{def:b1:c-strings-structs-files:text-binary}{text file} each digit is a character code and each line ends
with a newline (\code{0a}); its size depends on how many digits the numbers have. In the \omterm{def:b1:c-strings-structs-files:text-binary}{binary file} each number is
four bytes, whatever its value, least significant byte first.}
\end{omfigure}

The \omterm{def:b1:c-strings-structs-files:text-binary}{binary file} stores \code{1000 = 0x000003e8} as \code{e8 03 00 00}: the least significant byte first.

\begin{definition}[Byte order]\label{def:b1:c-strings-structs-files:byte-order}
The \emph{byte order}\index{byte order} (or endianness) of a machine is the order in which it stores the bytes of a
multi-byte integer at increasing addresses. A \emph{little-endian}\index{little-endian} machine stores the least
significant byte first, a \emph{big-endian}\index{big-endian} one the most significant byte first. x86-64 processors
are little-endian.
\end{definition}

\omcode{code/bachelor-1/06-c-strings-structs-files/listings/endian.c}{2}{18}{The bytes of an \code{int}, and a byte
widened two ways}\label{lst:b1:c-strings-structs-files:endian}

\omoutput{code/bachelor-1/06-c-strings-structs-files/outputs/endian.txt}{\omterm{def:b1:c-strings-structs-files:byte-order}{Little-endian}; and \code{0xf0} becomes $-16$
or 240.}\label{out:b1:c-strings-structs-files:endian}

\begin{definition}[Sign extension]\label{def:b1:c-strings-structs-files:sign-extension}
\emph{Sign extension}\index{sign extension} is the widening of a signed integer to more bits by copying its sign bit
into all the new high bits, which keeps its value in two's complement: the byte \code{f0} read as a signed
\code{char} ($-16$) becomes \code{fffffff0}, still $-16$. An unsigned value is widened by zeros instead.
\end{definition}

Whether plain \code{char} is signed depends on the platform (it is signed here); code that reads bytes from a file
uses \code{unsigned char}, as \cref{lst:b1:c-strings-structs-files:endian} does, or a byte above 127 becomes negative.
A \omterm{def:b1:c-strings-structs-files:text-binary}{binary file} written by \code{fwrite} from memory therefore depends on the machine twice: on its \omterm{def:b1:c-strings-structs-files:byte-order}{byte order} and on
the layout of its \omterm{def:b1:c-strings-structs-files:struct}{structures}. A format meant to be read elsewhere writes each value byte by byte, in an order the
format fixes (\cref{exo:b1:c-strings-structs-files:10}).

\begin{definition}[File position]\label{def:b1:c-strings-structs-files:position}
The \emph{file position}\index{file position} of a \omterm{def:b1:c-strings-structs-files:stream}{stream} is the offset, in bytes from the start of the file, at which
the next read or write happens; each read or write advances it by the number of bytes transferred. \code{ftell}
returns it, and \code{fseek} moves it, which allows access to any record of a file of fixed-size records without
reading the ones before.
\end{definition}

Reading byte after byte from the start is sequential access; jumping with \code{fseek} to offset $k \times s$, the
record number $k$ of size $s$, is direct access, possible because every record has the same size. A \omterm{def:b1:c-strings-structs-files:text-binary}{text file} of lines
of different lengths only allows the first.

\section{Reading input robustly}\label{sec:b1:c-strings-structs-files:robust}

Input comes from people and from other programs, and either can be wrong. \texttt{scanf("\%d", \&v)} skips spaces,
reads digits, and returns the number of values it converted; on a line \code{4.5} it reads 4 and leaves \code{.5} for
the next call, and on \code{twelve} it converts nothing and leaves the whole word in the input, so that a loop that
does not check its result reads the same word forever. The robust pattern reads whole lines, then converts them.

\begin{method}[Reading a number per line]\label{met:b1:c-strings-structs-files:robust}
\begin{enumerate}
\item Read a line with \code{fgets(buf, sizeof buf, f)}, which stops at the newline or at the buffer's size and returns
\code{NULL} at the end of the input.
\item Remove the newline: \texttt{buf[strcspn(buf, \string"\textbackslash{}n\string")] = '\textbackslash{}0';}.
\item Convert with \code{strtol} (or \code{strtod}), setting \code{errno} to 0 first; reject the line if no digit was
read, if characters remain after the number, or if \code{errno} is \code{ERANGE}.
\item Report a rejected line on \code{stderr}, with its number, and decide whether to go on or to stop.
\end{enumerate}
\end{method}

\omcode{code/bachelor-1/06-c-strings-structs-files/listings/read_numbers.c}{2}{27}{One number per line, every error
checked}\label{lst:b1:c-strings-structs-files:read}

\omoutput{code/bachelor-1/06-c-strings-structs-files/outputs/read_numbers.txt}{Nine lines of input, four
numbers: 12, 30, $-7$ and 100.}\label{out:b1:c-strings-structs-files:read}

The messages go to \code{stderr} so that they reach the person even when the standard output is redirected to a file
or to another program, which then receives only the result. The line \code{15 } (with a space) is rejected here; a
real program would probably accept trailing spaces, as \cref{exo:b1:c-strings-structs-files:12} does. With these
tools, the town's weather file can be read and checked.

\omcode{code/bachelor-1/06-c-strings-structs-files/listings/weather_csv.c}{2}{37}{Counting what the weather file
holds}\label{lst:b1:c-strings-structs-files:weather}

\omoutput{code/bachelor-1/06-c-strings-structs-files/outputs/weather.txt}{731 rows, two kinds of bad temperatures, and
a file that does not exist.}\label{out:b1:c-strings-structs-files:weather}

Two stations over 365 days make 730 rows: one row too many, which the weekend problem finds. Splitting a line with
\code{strtok}, the library's tokeniser, would have been shorter and wrong: it treats two consecutive commas as one, so
an empty field disappears and the next field takes its place.

\begin{safety}
Never use \code{gets}, removed from the language because it cannot be told the size of its buffer, nor \texttt{scanf("\%s")}
without a width. Copy strings with a size: \code{snprintf}, or a function that takes the capacity of the destination
and reports truncation. Check the result of every \code{fopen}, \code{fgets}, \code{fread}, \code{fwrite} and
\code{fclose}.
\end{safety}

\section{Exercises}\label{sec:b1:c-strings-structs-files:exercises}

\begin{exercise}[$\star$]\label{exo:b1:c-strings-structs-files:1}
Predict the eight numbers printed.
\omcode{code/bachelor-1/06-c-strings-structs-files/listings/exo_lengths.c}{5}{16}{Lengths and sizes}
\end{exercise}

\begin{exercise}[$\star$]\label{exo:b1:c-strings-structs-files:2}
Predict the sizes of the two \omterm{def:b1:c-strings-structs-files:struct}{structures}, drawing the offset of each field.
\omcode{code/bachelor-1/06-c-strings-structs-files/listings/exo_sizes.c}{2}{22}{Same \omterm{def:b1:c-strings-structs-files:struct}{fields}, two orders}
\end{exercise}

\begin{exercise}[$\star$]\label{exo:b1:c-strings-structs-files:3}
The program below reads its standard input, made of the four lines \code{1 12}, \code{2 0}, \code{3 x7} and \code{4 5}.
What does it print, and why is day 4 never read?
\omcode{code/bachelor-1/06-c-strings-structs-files/listings/exo_fscanf.c}{2}{12}{What \code{scanf} returns}
\end{exercise}

\begin{exercise}[$\star$]\label{exo:b1:c-strings-structs-files:4}
The program writes \code{0x12345678} with \code{fwrite}. Which four bytes does the file hold on this machine, in order,
and which on a \omterm{def:b1:c-strings-structs-files:byte-order}{big-endian} machine?
\omcode{code/bachelor-1/06-c-strings-structs-files/listings/exo_fwrite.c}{5}{14}{An \code{int} in a file}
\end{exercise}

\begin{exercise}[$\star\star$]\label{exo:b1:c-strings-structs-files:5}
Explain the report: which bytes does the first \code{strcat} write, and why is the write 13 bytes long? Give two
corrections.
\omcode{code/bachelor-1/06-c-strings-structs-files/listings/exo_strcat.c}{2}{12}{Building a greeting}
\omoutput[3-4,6-6,10-10]{code/bachelor-1/06-c-strings-structs-files/outputs/exo_strcat.txt}{Lines of the report.}
\end{exercise}

\begin{exercise}[$\star\star$]\label{exo:b1:c-strings-structs-files:6}
With the input \code{10}, \code{20}, \code{30} on three lines, this program prints \code{4 numbers, total 90}. Explain,
and correct the loop.
\omcode{code/bachelor-1/06-c-strings-structs-files/listings/exo_feof.c}{3}{15}{Until the end of the file}
\end{exercise}

\begin{exercise}[$\star\star$]\label{exo:b1:c-strings-structs-files:7}
Write \code{int str_append(char *dst, size_t cap, const char *src)}, which appends as much of \code{src} to the string
in \code{dst} (a buffer of \code{cap} bytes) as fits, always leaves \code{dst} terminated, and returns 0, or $-1$ if it
had to truncate.
\end{exercise}

\begin{exercise}[$\star\star$]\label{exo:b1:c-strings-structs-files:8}
A settings file holds lines \code{key=value}, empty lines and comment lines starting with \texttt{\#}. Write a function
that reads it into an \omterm{def:b1:pointers-and-arrays:array}{array} of \code{struct setting { char key[32]; char value[64]; }}, and reports the number of the
first malformed line.
\end{exercise}

\begin{exercise}[$\star\star$]\label{exo:b1:c-strings-structs-files:9}
The \omterm{def:b1:c-strings-structs-files:alignment}{alignments} are powers of two, and the size of each type is a multiple of its \omterm{def:b1:c-strings-structs-files:alignment}{alignment}. Prove that ordering the
\omterm{def:b1:c-strings-structs-files:struct}{fields} of a \omterm{def:b1:c-strings-structs-files:struct}{structure} by decreasing \omterm{def:b1:c-strings-structs-files:alignment}{alignment} gives the smallest possible size, namely the sum of the sizes of the
\omterm{def:b1:c-strings-structs-files:struct}{fields} rounded up to a multiple of the largest \omterm{def:b1:c-strings-structs-files:alignment}{alignment}.
\end{exercise}

\begin{exercise}[$\star\star\star$]\label{exo:b1:c-strings-structs-files:10}
Write \code{record_write} and \code{record_read} for the record \code{{ uint32_t date; int16_t temp; uint16_t
rain; }} in an 8-byte format, \omterm{def:b1:c-strings-structs-files:byte-order}{little-endian}, written byte by byte with shifts so that the file is the same on every
machine. Test the round trip, and give the 8 bytes for the date 20250614, $-35$ and 123.
\end{exercise}

\begin{exercise}[$\star\star\star$]\label{exo:b1:c-strings-structs-files:11}
A sentence of $n$ words of $w$ letters, each followed by a space, is built in a large buffer by $n$ calls of
\code{strcat}. Counting the bytes of the destination that \code{strcat} examines to find its end, then the bytes it
writes, show that the total is $(w+1)\,n(n-1)/2 + n(w+3)$. Give a linear method.
\end{exercise}

\begin{exercise}[$\star\star\star$]\label{exo:b1:c-strings-structs-files:12}
Write \code{int read_int(const char *s, int *out)}, which accepts optional spaces, a sign, digits and optional
trailing spaces or newline, and rejects an empty string, a non-number, trailing characters and any value outside the
range of \code{int}. Test it on a table of at least twelve cases.
\end{exercise}

\section{Problem: A binary file for the weather station}\label{sec:b1:c-strings-structs-files:problem}

\begin{problem}[{Weekend problem --- A binary file for the weather station}]\label{pb:b1:c-strings-structs-files:1}
The town's weather file holds one line per station and per day: date, station, temperature in degrees Celsius, rain in
millimetres and wind in kilometres per hour. A program reads it into memory, skipping a line identical to the one
before, and writes the readings to three archives: \code{first.bin} with the \omterm{def:b1:c-strings-structs-files:struct}{structure} \code{reading_first},
\code{reading.bin} with \code{reading}, and \code{packed.bin} in a format of 12 bytes written byte by byte. It then
reads the packed file back.
\omcode{code/bachelor-1/06-c-strings-structs-files/listings/weather.h}{2}{35}{Two \omterm{def:b1:c-strings-structs-files:struct}{structures} for one day of one
station}\label{lst:b1:c-strings-structs-files:weather-h}
\omcode{code/bachelor-1/06-c-strings-structs-files/listings/weather_io.c}{16}{48}{Parsing a
line}\label{lst:b1:c-strings-structs-files:weather-io}
\omcode{code/bachelor-1/06-c-strings-structs-files/listings/weather_io.c}{85}{93}{Packing a reading in 12
bytes}\label{lst:b1:c-strings-structs-files:pack}
\omoutput{code/bachelor-1/06-c-strings-structs-files/outputs/problem.txt}{The run, the sizes of the four files, and the
first bytes of two archives.}\label{out:b1:c-strings-structs-files:problem}

\needspace{4\baselineskip}
\medskip\noindent\textbf{Part I --- Designing the record.}
\begin{enumerate}
\item Give the offset of each field of \code{reading_first}, and the number of \omterm{def:b1:c-strings-structs-files:alignment}{padding} bytes after \code{station},
after \code{flag} and at the end. Check the total, 32.
\item Same question for \code{reading}, and check 24.
\item How many bytes of the 32 hold data? What fraction of \code{first.bin} is \omterm{def:b1:c-strings-structs-files:alignment}{padding}?
\item Why is \code{reading} not 22 bytes, the sum of its \omterm{def:b1:c-strings-structs-files:struct}{fields}? What would go wrong in an \omterm{def:b1:pointers-and-arrays:array}{array} of 22-byte \omterm{def:b1:c-strings-structs-files:struct}{structures}?
\item The temperature is stored as a \code{double} in memory but as a 2-byte count of tenths in the packed format. Which
temperatures can 2 signed bytes hold, and is that enough for the town?
\item The flag has two bits, \code{TEMP_MISSING} and \code{TEMP_ERROR}. Why store a flag rather than a special
temperature such as $-99.0$, the value the logger itself writes?
\end{enumerate}

\needspace{4\baselineskip}
\medskip\noindent\textbf{Part II --- Text against binary.}
\begin{enumerate}[resume]
\item Recompute the sizes of the three archives from the number of readings and the three record sizes.
\item The CSV file is 24\,358 bytes. Its header line has 36 characters. What is the mean length of a data line,
newline included?
\item In the first 32 bytes of \code{first.bin}, find the date, the station, the temperature 3.9 and the wind 15.3.
Decode the date \code{f5 fd 34 01}.
\item Bytes 17 to 19 of \code{first.bin} are \code{40 55 55}, not zeros. What are they, and why could writing \omterm{def:b1:c-strings-structs-files:alignment}{padding}
bytes to a file be a problem for a program that handles secrets?
\item Why are the \omterm{def:b1:c-strings-structs-files:alignment}{padding} bytes of \code{reading.bin} zeros? (Look at where \code{days} is declared in
\code{archive.c}.)
\item The program jumps straight to reading 160 of \code{packed.bin}. At which \omterm{def:b1:c-strings-structs-files:position}{file position} does it read, and why
could it not do the same in the CSV file?
\end{enumerate}

\needspace{4\baselineskip}
\medskip\noindent\textbf{Part III --- Byte order and a portable format.}
\begin{enumerate}[resume]
\item Decode the first 12 bytes of \code{packed.bin}, and the second reading (the next 12 bytes).
\item Which bytes would a \omterm{def:b1:c-strings-structs-files:byte-order}{big-endian} machine write for the first reading into \code{reading.bin} with \code{fwrite}? And
into \code{packed.bin}?
\item Why does \code{pack_reading} convert the temperature to \code{int16_t} before \code{uint16_t}? Give the two bytes
written for $-3.5$.
\item In \code{unpack_reading}, explain how \code{(int16_t)get(in + 6, 2)} recovers $-35$ from those two bytes, and
where \omterm{def:b1:c-strings-structs-files:sign-extension}{sign extension} happens.
\item The packed format rounds to tenths with \code{lround}. Why is that exact for this file, but would not be for a
temperature measured to the hundredth?
\item A reader on another machine receives \code{packed.bin}. What must the format's documentation say for it to read
the file correctly?
\end{enumerate}

\needspace{4\baselineskip}
\medskip\noindent\textbf{Part IV --- Reading the CSV robustly.}
\begin{enumerate}[resume]
\item The program reports 731 lines, 1 duplicate and 730 readings. Why is 730 the expected count, and how does
\code{read_weather} find the duplicate? What kind of duplicate would it miss?
\item Six temperatures are missing and two are error codes. How does \code{parse_reading} tell the three cases apart
(missing, $-99.0$, a real value)?
\item Why does \code{field_double} compare \code{end} with \code{stop}? Give a field it rejects that \code{strtod}
alone would accept.
\item What does \code{parse_reading} return for the line \code{2025-03-01,Hilltop,4.1,0.0} (a field missing), and
which test detects it?
\item The packed archive is how many times smaller than the CSV? Than \code{first.bin}?
\item State the named result of the problem, \emph{the size of the archive}.
\end{enumerate}
\end{problem}
